\documentclass[12pt]{article}
\usepackage{amsmath,amssymb,amsthm}
\usepackage[margin=1in]{geometry}
\usepackage{setspace}
\onehalfspacing

\newtheorem{proposition}{Proposition}

\title{Two-Period Example and Envelope Theorem \\ \medskip \large Following Fagereng, Gomez, Gouin-Bonenfant, Holm, Moll, and Natvik (2025)}
\author{}
\date{}

\begin{document}
\maketitle

\section{Setup}

Consider a two-period economy with $t = 0, 1$. Individual $i$ receives exogenous labor income $Y_{i,0}$ at $t=0$ and $Y_{i,1}$ at $t=1$. There is one asset with price $P_0 > 0$ at $t=0$ that pays a dividend $D_1 > 0$ at $t=1$.

Individual $i$ enters period $0$ with an endowment of $N_{i,-1}$ shares and chooses holdings $N_{i,0}$ to maximize lifetime utility:
\begin{equation}
  V_{i,0} = \max_{N_{i,0}} \; U(C_{i,0}) + \beta \, U(C_{i,1})
\end{equation}
subject to the budget constraints
\begin{align}
  C_{i,0} &= Y_{i,0} + N_{i,-1} P_0 - N_{i,0} P_0, \label{eq:bc0} \\
  C_{i,1} &= Y_{i,1} + N_{i,0} D_1, \label{eq:bc1}
\end{align}
where $U(\cdot)$ is increasing and strictly concave, and $\beta < 1$ is the discount factor.

The welfare effect of a change in asset prices is measured in money-metric terms as the \emph{equivalent variation} scaled by marginal utility:
\begin{equation}
  \frac{dV_{i,0}}{U'(C_{i,0})}.
\end{equation}
This expresses the welfare change in units of period-0 consumption.

\section{Envelope Theorem Result}

\begin{proposition}
  Applying the envelope theorem to the individual's problem, the welfare effect of a change in the asset price $P_0$ is:
  \begin{equation} \label{eq:envelope}
    \frac{dV_{i,0}}{U'(C_{i,0})} = (N_{i,-1} - N_{i,0}) \, dP_0.
  \end{equation}
\end{proposition}

\noindent \textbf{Interpretation.} What matters for welfare is \emph{net asset sales} $(N_{i,-1} - N_{i,0})$, not asset holdings. An individual who holds assets but does not trade ($N_{i,0} = N_{i,-1}$) experiences zero welfare effect from a price change --- paper gains and losses are irrelevant for non-traders. Only individuals who are net sellers ($N_{i,-1} > N_{i,0}$) benefit from a price increase, and net buyers ($N_{i,0} > N_{i,-1}$) are hurt.

\section{Proof via the Lagrangian}

Write the Lagrangian for the individual's problem, embedding the budget constraints \eqref{eq:bc0}--\eqref{eq:bc1}:
\begin{equation}
  \mathcal{L} = U\bigl(Y_{i,0} + N_{i,-1} P_0 - N_{i,0} P_0\bigr) + \beta \, U\bigl(Y_{i,1} + N_{i,0} D_1\bigr).
\end{equation}

\noindent \textbf{Step 1: First-order condition.} Differentiating with respect to the choice variable $N_{i,0}$:
\begin{equation} \label{eq:foc}
  \frac{\partial \mathcal{L}}{\partial N_{i,0}} = -P_0 \, U'(C_{i,0}) + \beta \, D_1 \, U'(C_{i,1}) = 0.
\end{equation}
This is the standard Euler equation: $P_0 \, U'(C_{i,0}) = \beta \, D_1 \, U'(C_{i,1})$.

\noindent \textbf{Step 2: Total differentiation with respect to $P_0$.} Totally differentiate the value function (at the optimum) with respect to $P_0$. By the envelope theorem, we can ignore the indirect effect through the optimal choice $N_{i,0}^*$ (since $\partial \mathcal{L}/\partial N_{i,0} = 0$ at the optimum). This gives:
\begin{equation}
  \frac{dV_{i,0}}{dP_0} = \frac{\partial \mathcal{L}}{\partial P_0} = (N_{i,-1} - N_{i,0}) \, U'(C_{i,0}).
\end{equation}
To see this, differentiate the Lagrangian with respect to $P_0$ holding $N_{i,0}$ fixed:
\begin{equation}
  \frac{\partial \mathcal{L}}{\partial P_0} = U'(C_{i,0}) \cdot (N_{i,-1} - N_{i,0}).
\end{equation}
The term $(N_{i,-1} - N_{i,0})$ arises because $P_0$ appears in the period-0 budget constraint \eqref{eq:bc0} multiplied by both $N_{i,-1}$ (asset endowment value) and $-N_{i,0}$ (cost of new purchases). The period-1 constraint \eqref{eq:bc1} does not involve $P_0$, so there is no direct effect through $C_{i,1}$.

\noindent \textbf{Step 3: Money-metric welfare.} Dividing both sides by $U'(C_{i,0})$:
\begin{equation}
  \frac{dV_{i,0}}{U'(C_{i,0})} = (N_{i,-1} - N_{i,0}) \, dP_0,
\end{equation}
which is the desired result \eqref{eq:envelope}. \hfill $\blacksquare$

\end{document}
